% Experiments --- setup, datasets, metrics, harness. Numbers live in
% Results; this section fixes the protocol. The LLM-benchmark table is a
% honest placeholder pending US-069 (H100 eval), per the data-real rule.
\section{Configuración Experimental}
\label{sec:experiments}

\subsection{Conjuntos de datos}
\paragraph{PASTIS-R}~\cite{garnot2021pastis} es el benchmark principal:
2{,}433 parches Sentinel-2 multitemporales ($128\times128$ px), 18 clases
de cultivo, Francia metropolitana, con cinco folds espacialmente
disjuntos. Salvo que se indique lo contrario, los modelos de
segmentación se evalúan en un fold de prueba retenido y los combinadores
de ensamble sobre predicciones out-of-fold.

\paragraph{Datasets tipo PASTIS construidos a mano (Italia, Alemania).}
Para medir la transferencia con datos reales donde no existe un benchmark en
formato PASTIS, construimos dos datasets a mano en el mismo formato que PASTIS-R.
\emph{Italia (Toscana/Pianura, 2018)}: $1{,}438$ parches de $128\times128$\,px y
$103{,}350$ parcelas etiquetadas. \emph{Alemania (Baja Sajonia, DE4, 2023)}:
$1{,}246$ parches y $53{,}460$ parcelas con una ventana de 14 meses ($\approx156$k
parcelas etiquetadas en total). Ambos siguen la misma receta: polígonos con el
cultivo declarado de EuroCrops como etiquetas, la serie temporal Sentinel-2
completa por parche (24 fechas en Italia, ~41 en Alemania) con una máscara de
cultivo píxel a píxel, un resumen AlphaEarth por parcela vía Google Earth Engine,
y una división en cinco bloques geográficos para una evaluación sin fugas.
Construir estos datasets es lo que hace medible la transferencia
Francia$\rightarrow$Italia y Francia$\rightarrow$Alemania; se versionan con DVC.

\paragraph{Sen4AgriNet (Cataluña)}~\cite{sykas2022sen4agrinet}
(CC-BY-SA-4.0) se usa para el estudio de transferencia
Francia$\rightarrow$Cataluña. Un subconjunto de Cataluña se versiona con
DVC (\texttt{data/sen4agrinet.dvc}).

\paragraph{EuroCropsML}~\cite{reuss2025eurocropsml} (CC-BY-SA-4.0)
sustenta la curva few-shot de $k$-shots; el subconjunto crudo se versiona
con DVC (\texttt{data/transfer/eurocropsml.dvc}).

\paragraph{AgroMind}~\cite{ruan2025agromind} se usa solo para evaluar la
capa conversacional. No tiene partición de entrenamiento (alrededor de
28{,}482 pares de QA); el ajuste fino sobre él induciría fuga, por lo que
es estrictamente un conjunto de evaluación.

\subsection{Métricas}
La segmentación se reporta con la intersección sobre la unión media
(mIoU) y el F1 promediado en macro a nivel de parcela; también
reportamos la exactitud de píxel donde sea relevante. El promedio en
macro es el objetivo principal porque penaliza por igual los errores en
cultivos raros y comunes, lo que coincide con el objetivo de negocio.
Los combinadores de ensamble se comparan por F1 macro y exactitud
out-of-fold.

\subsection{Arnés de evaluación y reproducibilidad}
Todas las corridas de segmentación se puntúan con un arnés de evaluación
compartido (\texttt{ml/eval/}) que impone folds espacialmente disjuntos
y aborta ante un solapamiento geográfico entre los conjuntos de
entrenamiento y evaluación. Un módulo de remapeo de etiquetas
(\texttt{ml/eval/class\_remap.py}) proyecta las predicciones sobre el
espacio macro de etiquetas armonizado para la comparación entre
regiones. Cada cifra reportada en la Sección~\ref{sec:results} está
anclada a un artefacto versionado (JSON/CSV/registro) en el repositorio,
registrado en comentarios \texttt{\% src:} en el código fuente; los
conjuntos de datos y los pesos se rastrean con DVC y las corridas con
MLflow.

\subsection{Protocolo de benchmark de la capa conversacional}
\label{sec:experiments-llm}
Evaluamos la capa conversacional en dos ejes independientes. El primero
es un eje de \emph{ingeniería} --- la intercambiabilidad de backends ---
que ya está medido: detrás de un perceptor fijo, se envía una misma
consulta anclada a los tres backends de servicio efectivamente
desplegados (Gemini~3.5~Flash en la nube, más Qwen3~MoE y
Qwen~3.6-VL on-premises), y se registra si la respuesta es idéntica y su
latencia de extremo a extremo. El resultado medido se reporta en la
Sección~\ref{sec:results} (Tabla~\ref{tab:backend-switch}). El segundo es
un eje de \emph{calidad} --- un protocolo LLM-como-juez sobre AgroMind que
puntúa las respuestas del razonador congelado por exactitud y anclaje ---
que depende de una corrida de evaluación separada en la GPU H100 (US-069)
y sigue pendiente; reportamos el protocolo y diferimos las cifras en
lugar de fabricarlas. No se reporta ninguna puntuación sintética de
AgroMind en ninguna parte de este artículo.
